% Wersja źródłowa: audyt II, jednostka 17
\setcounter{section}{16}
\section{Przestrzenie funkcji gładkich i przestrzeń Schwartza}
\label{sec:gladkie-schwartz17}

\subsection{Topologia przestrzeni \texorpdfstring{$C^\infty(\Omega)$}{Cinfty(Omega)}}

Niech $d\ge1$, niech $\Omega\subseteq\R^d$ będzie niepustym zbiorem
otwartym, a $\K\in\{\R,\C\}$. Dla wielowskaźnika
$\alpha=(\alpha_1,\dots,\alpha_d)\in\N_0^d$ oznaczamy
$|\alpha|=\sum_j\alpha_j$, $\alpha!=\prod_j\alpha_j!$,
$x^\alpha=\prod_jx_j^{\alpha_j}$ oraz
$\partial^\alpha=\partial_1^{\alpha_1}\cdots\partial_d^{\alpha_d}$.
Dla $\gamma,\beta\in\N_0^d$ piszemy $\gamma\le\beta$,
gdy $\gamma_j\le\beta_j$ dla każdego $j$;
wtedy $\binom{\beta}{\gamma}=\prod_{j=1}^d\binom{\beta_j}{\gamma_j}$.
Symbole $e_i$ oznaczają wektory standardowej bazy $\R^d$,
traktowane także jako wielowskaźniki. Piszemy $K\Subset\Omega$,
gdy $K$ jest zwartym podzbiorem $\Omega$.
Przez $C^\infty(\Omega;\K)$ rozumiemy przestrzeń funkcji
$\Omega\to\K$ posiadających ciągłe pochodne cząstkowe wszystkich rzędów.
Przez $|x|$ dla $x\in\R^d$ rozumiemy normę euklidesową,
a przez $\lambda_d$ miarę Lebesgue'a na $\R^d$.
Nie narzucamy ograniczoności funkcji ani ich pochodnych na całym $\Omega$.

\begin{definicja}[Seminormy zbieżności lokalnej wszystkich pochodnych]
\label{def:seminormy-cinf17}
Niech $d\ge1$, $\Omega\subseteq\R^d$ będzie niepustym zbiorem
otwartym, a $\K=\R$ lub $\K=\C$.
Dla zwartego $K\subseteq\Omega$ i liczby $m\in\N_0$ określamy
\[
 p_{K,m}(f)=\max_{|\alpha|\le m}\sup_{x\in K}|\partial^\alpha f(x)|,
 \qquad f\in C^\infty(\Omega;\K).
\]
Dla $K=\varnothing$ przyjmujemy $p_{K,m}=0$.
Nierówność trójkąta i jednorodność wynikają z odpowiednich
własności wartości bezwzględnej, supremum i maksimum.
Topologią standardową na $C^\infty(\Omega;\K)$ nazywamy
topologię generowaną przez rodzinę tych seminorm.
Sieć $(f_\iota)$ zbiega do $f$ w tej topologii wtedy i tylko wtedy,
gdy dla każdego zwartego $K\subseteq\Omega$ i każdego
$\alpha\in\N_0^d$ pochodne $\partial^\alpha f_\iota$
zbiegają jednostajnie do $\partial^\alpha f$ na $K$.
\end{definicja}

\begin{lemat}[Wyczerpanie zbiorami zwartymi]
\label{lem:exhaustion17}
Dla każdego niepustego otwartego $\Omega\subseteq\R^d$ istnieje
ciąg zbiorów zwartych $(K_n)_{n\ge1}$ taki, że
$K_n\subseteq\operatorname{int}K_{n+1}$, $\bigcup_nK_n=\Omega$
i każdy zwarty $K\subseteq\Omega$ jest zawarty w pewnym $K_n$.
\end{lemat}
\begin{proof}
Dla $\Omega\ne\R^d$ weźmy
$K_n=\{x\in\Omega:|x|\le n,\
\operatorname{dist}(x,\R^d\setminus\Omega)\ge1/n\}$;
dla $\Omega=\R^d$ przyjmijmy $K_n=\{x:|x|\le n\}$.
Zbiory te są zwarte: warunek odległości jest domknięty,
a dodatnia odległość od dopełnienia wymusza przynależność do $\Omega$.
Nierówności definiujące $K_n$ są ściślejsze niż analogiczne dla
$K_{n+1}$, stąd $K_n\subseteq\operatorname{int}K_{n+1}$.
Każdy $x\in\Omega$ ma dodatnią odległość od dopełnienia
i skończoną normę, więc należy do pewnego $K_n$.
Jeżeli zwarty $K\subseteq\Omega$ nie byłby zawarty w żadnym $K_n$,
można byłoby wybrać $x_n\in K\setminus K_n$.
Ze zwartości $K$ istnieje podciąg $(x_{n_j})$ zbieżny
do $x\in K$. Pewne $K_N$ zawiera $x$ w swoim wnętrzu,
bo $x\in K_{N-1}\subseteq\operatorname{int}K_N$
dla dostatecznie dużego $N$.
Dla dużych $j$ mamy $n_j\ge N$ i $x_{n_j}\in K_N\subseteq K_{n_j}$,
sprzeczność. Gdy dopełnienie jest puste, argument jest identyczny.
\end{proof}

\begin{twierdzenie}[Przeliczalna rodzina seminorm]
\label{thm:countable-cinf17}
Niech $\Omega\subseteq\R^d$ będzie niepusty i otwarty,
a $(K_n)$ ciągiem z Lematu~\ref{lem:exhaustion17}.
Wówczas topologia z Definicji~\ref{def:seminormy-cinf17}
jest generowana przez rosnący ciąg seminorm
$q_n=p_{K_n,n}$, $n\ge1$ (w szczególności $q_n=0$,
jeżeli $K_n=\varnothing$). Jest ona metryzowalna i Hausdorffa.
Ciąg $(f_j)$ zbiega w niej do $f$ wtedy i tylko wtedy, gdy
$q_n(f_j-f)\to0$ dla każdego $n\ge1$.
\end{twierdzenie}
\begin{proof}
Każde $q_n$ jest jedną z seminorm $p_{K,m}$.
Z drugiej strony dla ustalonych $K\Subset\Omega$ i $m$
wybierzmy $n\ge m$ z $K\subseteq K_n$.
Wtedy $p_{K,m}\le q_n$, więc obie rodziny wyznaczają tę samą
topologię. Ponieważ $K_n\subseteq K_{n+1}$, seminormy $q_n$
są niemalejące. Jeśli $f\ne0$, to $f(x)\ne0$ dla pewnego
$x\in\Omega$ i $x\in K_n$ dla pewnego $n$, zatem $q_n(f)>0$.
Topologia jest więc Hausdorffa. Jest metryzowalna na mocy
kryterium dla przeliczalnej rodziny seminorm; można użyć metryki
$\sum_{n\ge1}2^{-n}\min\{1,q_n(f-g)\}$.
Z definicji topologii wyznaczonej przez seminormy wynika
podane kryterium zbieżności.
\end{proof}

\begin{twierdzenie}[Zupełność $C^\infty(\Omega)$]
\label{thm:frechet-cinf17}
Dla każdego niepustego otwartego $\Omega\subseteq\R^d$ przestrzeń
$C^\infty(\Omega;\K)$ ze standardową topologią jest przestrzenią
Frécheta.
\end{twierdzenie}
\begin{proof}
Metryzowalność, lokalna wypukłość i warunek Hausdorffa wynikają
z Twierdzenia~\ref{thm:countable-cinf17}; pozostaje zupełność.
Niech $(f_j)$ będzie ciągiem Cauchy'ego dla wszystkich $q_n$.
Dla każdego wielowskaźnika $\alpha$ pochodne
$\partial^\alpha f_j$ są jednostajnie Cauchy'ego na każdym
$K_n$, gdy $n\ge|\alpha|$. Z zupełności przestrzeni
$C(K_n;\K)$ mają granice jednostajne $g_{\alpha,n}$.
Na przecięciach niepustych $K_n$ granice się zgadzają
na mocy jednoznaczności granicy punktowej. Określają więc
funkcje $g_\alpha:\Omega\to\K$, ciągłe, ponieważ każdy punkt
$\Omega$ leży we wnętrzu pewnego $K_n$ z $n\ge|\alpha|$; w szczególności
$f_j\to g_0$ lokalnie jednostajnie.

Wykażemy, że $g_\alpha=\partial^\alpha g_0$.
Ustalmy $x\in\Omega$ i zamkniętą kulę $\overline B(x,r)\Subset\Omega$.
Dla $i\in\{1,\dots,d\}$ i $|t|<r$, z twierdzenia
Newtona--Leibniza na odcinku od $x$ do $x+te_i$ zachodzi
\[
 \partial^\alpha f_j(x+te_i)-\partial^\alpha f_j(x)
 =\int_0^t\partial^{\alpha+e_i}f_j(x+se_i)\,\mathrm ds.
\]
Wybierając $K_n$ zawierający tę kulę i $n\ge|\alpha|+1$,
możemy przejść do granicy po obu stronach: błąd całki
nie przekracza $|t|\sup_{K_n}|\partial^{\alpha+e_i}f_j-
g_{\alpha+e_i}|$, który dąży do zera. Dla $t\ne0$ mamy zatem
\[
\frac{g_\alpha(x+te_i)-g_\alpha(x)}{t}
=\frac1t\int_0^t g_{\alpha+e_i}(x+se_i)\,\mathrm ds.
\]
Z ciągłości $g_{\alpha+e_i}$ prawa strona dąży do
$g_{\alpha+e_i}(x)$ przy $t\to0$, również dla $t<0$.
Stąd $\partial_i g_\alpha=g_{\alpha+e_i}$ na $\Omega$.
Indukcja względem $|\alpha|$ daje $g_0\in C^\infty(\Omega)$
i wszystkie wymagane tożsamości. Jednostajna zbieżność
pochodnych na $K_n$ oznacza $q_n(f_j-g_0)\to0$ dla każdego $n$.
Ciąg jest Cauchy'ego dla metryki z poprzedniego dowodu
wtedy i tylko wtedy, gdy jest Cauchy'ego dla każdej $q_n$:
w jedną stronę wystarczy rozważyć $n$-ty składnik sumy,
w drugą najpierw oszacować ogon szeregu przez $\sum_{n>N}2^{-n}$,
a następnie skończoną sumę dla $n\le N$.
Wykazaliśmy zatem zupełność tej metryki zgodnej z topologią
przestrzeni i niezmienniczej względem przesunięć.
\end{proof}

\begin{stwierdzenie}[Operacje i ograniczoność]
\label{prop:operations-cinf17}
Niech $\Omega\subseteq\R^d$ będzie niepusty i otwarty.
W przestrzeni $C^\infty(\Omega;\K)$ różniczkowanie
$\partial^\alpha$ jest operatorem liniowym ciągłym,
a mnożenie $(f,g)\mapsto fg$ jest odwzorowaniem dwuliniowym
łącznie ciągłym. Zbiór $B\subseteq C^\infty(\Omega;\K)$
jest ograniczony w sensie topologii przestrzeni liniowej
topologicznej wtedy i tylko wtedy, gdy
$\sup_{f\in B}p_{K,m}(f)<\infty$ dla wszystkich zwartych
$K\Subset\Omega$ i wszystkich $m\in\N_0$.
\end{stwierdzenie}
\begin{proof}
Z nierówności $p_{K,m}(\partial^\alpha f)\le p_{K,m+|\alpha|}(f)$
wynika ciągłość różniczkowania. Ze wzoru Leibniza
$\partial^\beta(fg)=\sum_{\gamma\le\beta}\binom{\beta}{\gamma}
\partial^\gamma f\,\partial^{\beta-\gamma}g$
oraz $\sum_{\gamma\le\beta}\binom{\beta}{\gamma}=2^{|\beta|}$
wynika $p_{K,m}(fg)\le 2^m p_{K,m}(f)p_{K,m}(g)$,
co dowodzi łącznej ciągłości mnożenia w zerze;
dwuliniowość i ciągłość przy dowolnej parze wynikają przez
rozwinięcie $(f+h)(g+k)-fg$. Kryterium ograniczoności
wynika z ogólnego kryterium seminormowego dla przestrzeni
lokalnie wypukłych: rodzina obrazów zbioru ograniczonego przez
każdą ciągłą seminormę jest ograniczona w $\R$ i odwrotnie.
\end{proof}

\subsection{Funkcje odcięcia i przestrzeń dualna topologiczna}

\begin{lemat}[Gładka funkcja odcięcia]
\label{lem:cutoffs17}
Jeżeli $K\subseteq\Omega\subseteq\R^d$, gdzie $K$ jest zwarty,
a $\Omega$ otwarty, to istnieje $\chi\in C_c^\infty(\Omega;\R)$
taka, że $0\le\chi\le1$ i $\chi=1$ w pewnym otoczeniu $K$.
Tu $C_c^\infty$ oznacza funkcje gładkie o zwartym nośniku.
\end{lemat}
\begin{proof}
Jeżeli $K=\varnothing$, wystarczy funkcja zerowa.
Dla $K\ne\varnothing$ wybierzmy otwarte zbiory $V,W$ takie, że
$K\subseteq V$, $\overline V\subseteq W$, $\overline W\Subset\Omega$.
Dla każdego $a\in K$ wybieramy $r_a>0$ tak,
aby $\overline B(a,3r_a)\subseteq\Omega$. Ze zwartości $K$
wybieramy skończone pokrycie kulami $B(a_j,r_{a_j})$.
Jako $V$ bierzemy sumę tych kul, a jako $W$ sumę kul
$B(a_j,2r_{a_j})$. Domknięcia skończonych sum są sumami
domknięć, więc otrzymujemy żądane zawierania.
Odległości $\operatorname{dist}(\overline V,\R^d\setminus W)$
i $\operatorname{dist}(K,\R^d\setminus V)$ są dodatnie,
ponieważ $\overline V$ i $K$ są zwarte
i rozłączne z odpowiednimi zbiorami domkniętymi.
Wybieramy $\delta>0$ mniejszą od obu odległości, tak że zbiór
$\{x:\operatorname{dist}(x,\overline V)\le\delta\}$
leży w $W$ oraz $\operatorname{dist}(K,\R^d\setminus V)>\delta$.
Niech $\rho\in C_c^\infty(\R^d)$ będzie nieujemna,
niech $\operatorname{supp}\rho\subseteq B(0,\delta)$ i
$\int_{\R^d}\rho(x)\,\mathrm d\lambda_d(x)=1$.
Taką funkcję otrzymujemy przez unormowanie
$x\mapsto\exp[-1/(1-4|x|^2/\delta^2)]$ dla $|x|<\delta/2$,
przedłużonej zerem dla $|x|\ge\delta/2$.
Gładkość przedłużenia wynika z faktu, że każda pochodna
funkcji $e^{-1/t}$ dla $t>0$ jest wielomianem w $1/t$
pomnożonym przez $e^{-1/t}$, a więc dąży do zera dla $t\downarrow0$.
Nośnik leży w $\overline B(0,\delta/2)\subset B(0,\delta)$.
Definiujemy $\chi(x)=\int_{\R^d}\chi_V(y)\rho(x-y)\,\mathrm d\lambda_d(y)$,
gdzie $\chi_V$ jest funkcją charakterystyczną $V$.
Przy każdym różniczkowaniu ilorazy różnicowe pochodnych
$\rho(x-y)$ są ograniczone przez supremum odpowiedniej
pochodnej wyższego rzędu, na mocy twierdzenia o wartości średniej.
Po pomnożeniu przez $\chi_V$ otrzymujemy całkowalną majorantę,
gdyż $V$ jest ograniczony. Z twierdzenia o zbieżności
zmajoryzowanej, indukcyjnie,
$\partial^\alpha\chi(x)=\int_V\partial^\alpha\rho(x-y)\,\mathrm d\lambda_d(y)$;
ciągłość tych pochodnych wynika z tego samego twierdzenia.
Stąd $\chi\in C^\infty(\R^d;\R)$.
Po podstawieniu $z=x-y$ mamy
$\chi(x)=\int_{\R^d}\chi_V(x-z)\rho(z)\,\mathrm d\lambda_d(z)$.
Ponieważ $0\le\chi_V\le1$, $\rho\ge0$ i całka $\rho$ jest równa $1$,
otrzymujemy $0\le\chi(x)\le1$.
Jeżeli $\operatorname{dist}(x,K)<\delta/2$ i
$z\in\operatorname{supp}\rho$, to
$\operatorname{dist}(x-z,K)\le\operatorname{dist}(x,K)+|z|<\delta$.
Z wyboru $\delta$ wynika $x-z\in V$, więc w tym otoczeniu
$K$ całka jest równa $1$, czyli $\chi(x)=1$. Jeżeli $\operatorname{dist}(x,\overline V)>\delta/2$,
to $\rho(x-y)=0$ dla każdego $y\in V$.
Stąd $\operatorname{supp}\chi\subseteq
\{x:\operatorname{dist}(x,\overline V)\le\delta/2\}$;
jest to zbiór zwarty zawarty w $W\Subset\Omega$,
co kończy dowód.
\end{proof}

\begin{wniosek}[Gęstość funkcji o zwartym nośniku]
\label{cor:density-cinf17}
Dla każdego niepustego otwartego $\Omega\subseteq\R^d$
przestrzeń $C_c^\infty(\Omega;\K)$ jest gęsta w
$C^\infty(\Omega;\K)$ ze standardową topologią.
\end{wniosek}
\begin{proof}
Dla wyczerpania $(K_n)$ wybierzmy z Lematu~\ref{lem:cutoffs17}
$\chi_n\in C_c^\infty(\Omega)$ równą $1$ na otoczeniu $K_n$.
Dla każdego $f\in C^\infty(\Omega)$ funkcje $\chi_n f$
mają zwarty nośnik. Dla ustalonego $K_m$ i dowolnego
rzędu pochodnych mamy $\chi_n f=f$ na otoczeniu $K_m$
dla wszystkich $n\ge m$. Zatem $\chi_n f\to f$
w topologii $C^\infty(\Omega)$.
\end{proof}

\begin{twierdzenie}[Kryterium ciągłości funkcjonału]
\label{thm:dual-cinf17}
Niech $\Omega\subseteq\R^d$ będzie niepusty i otwarty,
a $T:C^\infty(\Omega;\K)\to\K$ liniowy.
Następujące warunki są równoważne:
\begin{enumerate}
\item $T$ jest ciągły;
\item istnieją zwarty $K\Subset\Omega$, $m\in\N_0$ i $C>0$ takie, że
$|T(f)|\le C\max_{|\alpha|\le m}\sup_{x\in K}|\partial^\alpha f(x)|$
dla każdego $f\in C^\infty(\Omega;\K)$.
\end{enumerate}
W szczególności ciągły funkcjonał $T$ zależy tylko od wartości
pochodnych rzędów co najwyżej $m$ na pewnym zwartym zbiorze $K$:
jeżeli $\partial^\alpha f|_K=\partial^\alpha g|_K$ dla każdego
$|\alpha|\le m$, to $T(f)=T(g)$.
W szczególności jeśli $f$ znika w otoczeniu $K$, to $T(f)=0$.
\end{twierdzenie}
\begin{proof}
Wybierzmy rosnącą rodzinę seminorm $q_n$ z
Twierdzenia~\ref{thm:countable-cinf17}.
Ciągłość $T$ w zerze oznacza, że pewne otoczenie
$\{f:q_n(f)<\varepsilon\}$ jest odwzorowywane w
$\{z:|z|<1\}$. Jeśli $q_n(f)>0$, to dla każdego $0<a<\varepsilon/q_n(f)$
mamy $q_n(af)<\varepsilon$, zatem $a|T(f)|<1$.
Przechodząc z $a$ do $\varepsilon/q_n(f)$ od dołu,
otrzymujemy $|T(f)|\le\varepsilon^{-1}q_n(f)$.
Jeśli $q_n(f)=0$, to $q_n(af)=0$ dla każdego $a>0$,
więc $a|T(f)|<1$ dla wszystkich $a>0$, co daje $T(f)=0$.
To dowodzi warunku 2 z $K=K_n$ i $m=n$.
W razie potrzeby zwiększamy $n$, aby $K_n$ był niepusty;
nierówność pozostaje prawdziwa, gdyż $q_n$ są niemalejące.
Odwrotna implikacja wynika z ciągłości seminormy $p_{K,m}$.
Jeżeli wskazane pochodne $f$ i $g$ zgadzają się na $K$,
to $p_{K,m}(f-g)=0$, więc z liniowości
$|T(f)-T(g)|=|T(f-g)|\le Cp_{K,m}(f-g)=0$.
Jeśli $f=0$ na otoczeniu $K$, wszystkie jej pochodne
znikają na $K$ i można przyjąć $g=0$.
\end{proof}

\begin{przyklad}[Ewaluacja i pochodne ewaluacji]
\label{ex:evaluation-cinf17}
Dla $a\in\Omega$ i $\alpha\in\N_0^d$ funkcjonały
$T_{a,\alpha}(f)=(-1)^{|\alpha|}\partial^\alpha f(a)$
są ciągłe na $C^\infty(\Omega;\K)$, ponieważ
$|T_{a,\alpha}(f)|\le p_{\{a\},|\alpha|}(f)$.
Dla $\alpha=0$
przyjmują postać ewaluacji $f\mapsto f(a)$.
W teorii dystrybucji pojawią się jako pochodne
dystrybucji Diraca. Obecnie są po prostu przykładami
elementów przestrzeni dualnej topologicznej.
\end{przyklad}

\subsection{Przestrzeń Schwartza \texorpdfstring{$\mathcal S(\R^d)$}{S(Rd)}}

\begin{definicja}[Przestrzeń Schwartza]
\label{def:schwartz17}
Dla $d\ge1$ i $\K\in\{\R,\C\}$ definiujemy
\[
\mathcal S(\R^d;\K)=\left\{f\in C^\infty(\R^d;\K):
\sup_{x\in\R^d}|x^\alpha\partial^\beta f(x)|<\infty
\ \text{dla wszystkich }\alpha,\beta\in\N_0^d\right\}.
\]
Jej standardową topologię określa rodzina seminorm
$s_{\alpha,\beta}(f)=\sup_x|x^\alpha\partial^\beta f(x)|$.
Równoważną rodzinę tworzą
\[
 Q_{N,m}(f)=\max_{|\beta|\le m}\sup_{x\in\R^d}
 (1+|x|)^N|\partial^\beta f(x)|,
 \qquad N,m\in\N_0.
\]
Przyjmujemy też rosnący ciąg $Q_n=Q_{n,n}$, $n\ge1$.
\end{definicja}

\begin{lemat}[Równoważność seminorm Schwartza]
\label{lem:schwartz-seminorms17}
Seminormy $s_{\alpha,\beta}$ oraz $Q_{N,m}$ generują tę samą
topologię. Ponadto $Q_n$ są niemalejące, oddzielają punkty
i wyznaczają topologię przestrzeni Schwartza.
\end{lemat}
\begin{proof}
Ponieważ $|x^\alpha|\le |x|^{|\alpha|}\le(1+|x|)^{|\alpha|}$,
mamy $s_{\alpha,\beta}\le Q_{|\alpha|,|\beta|}$.
Z drugiej strony nierówność
$(1+|x|)^N\le C_{d,N}\bigl(1+\sum_{|\alpha|=N}|x^\alpha|\bigr)$
dla $N\ge1$ wynika z oszacowania
$(1+r)^N\le 2^N\max\{1,r^N\}$ dla $r\ge0$ oraz
$|x|^N\le d^{N/2}\sum_{j=1}^d|x_j|^N$.
Dla $N=0$ wystarczy stała $1$. W szczególności dla $N\ge1$
\[
Q_{N,m}(f)\le 2^N\sum_{|\beta|\le m}
\left(s_{0,\beta}(f)+d^{N/2}\sum_{j=1}^d s_{Ne_j,\beta}(f)\right).
\]
Dla $N=0$ mamy $Q_{0,m}=\max_{|\beta|\le m}s_{0,\beta}$.
Obie rodziny seminorm wyznaczają więc tę samą topologię.
Rodzina $Q_n$ jest współkońcowa względem $Q_{N,m}$,
gdy $n\ge\max\{N,m\}$. Oddzielanie punktów wynika już z $Q_1$.
\end{proof}

\begin{twierdzenie}[Zupełność przestrzeni Schwartza]
\label{thm:schwartz-frechet17}
Dla $\K=\R$ lub $\K=\C$ przestrzeń $\mathcal S(\R^d;\K)$
ze standardową topologią jest przestrzenią Frécheta.
Naturalne włożenie $\mathcal S(\R^d;\K)\hookrightarrow
C^\infty(\R^d;\K)$ jest ciągłe.
\end{twierdzenie}
\begin{proof}
Metryzowalność, lokalna wypukłość i warunek Hausdorffa wynikają z
Lematu~\ref{lem:schwartz-seminorms17}.
Niech $(f_j)$ będzie ciągiem Cauchy'ego dla wszystkich $Q_n$.
Ponieważ na każdym zwartym zbiorze i dla każdego rzędu pochodnych
seminormy $C^\infty$ są kontrolowane przez odpowiednie $Q_n$,
ciąg jest Cauchy'ego w $C^\infty(\R^d)$.
Na mocy Twierdzenia~\ref{thm:frechet-cinf17} istnieje
$f\in C^\infty(\R^d)$, do której $f_j$ zbiega lokalnie
wraz ze wszystkimi pochodnymi.
Ustalmy $N,m$ i $\varepsilon>0$. Istnieje $j_0$ takie, że
$Q_{N,m}(f_j-f_k)\le\varepsilon$ dla $j,k\ge j_0$.
Dla ustalonego $j\ge j_0$, $x\in\R^d$ i $|\beta|\le m$
mamy $(1+|x|)^N|\partial^\beta f_j(x)-\partial^\beta f_k(x)|
\le\varepsilon$ dla każdego $k\ge j_0$.
Przechodząc z $k$ do nieskończoności przy ustalonym $x$,
a następnie biorąc supremum po $x$ i maksimum po $\beta$,
otrzymujemy $Q_{N,m}(f_j-f)\le\varepsilon$.
Ponieważ $Q_{N,m}(f_j)<\infty$, wynika stąd
$Q_{N,m}(f)<\infty$. Wszystkie takie seminormy są skończone,
więc $f\in\mathcal S$, a $f_j\to f$ w topologii Schwartza.
Ciągłość włożenia wynika z nierówności
$p_{K,m}(f)\le Q_{0,m}(f)$ dla każdego $K\Subset\R^d$.
\end{proof}

\begin{przyklad}[Funkcje w przestrzeni Schwartza]
\label{ex:schwartz-functions17}
Niech $d\ge1$ i $\K\in\{\R,\C\}$.
Każda funkcja z $C_c^\infty(\R^d;\K)$ należy do
$\mathcal S(\R^d;\K)$, ponieważ wszystkie jej pochodne
znikają poza jej zwartym nośnikiem, a na nim
ważone pochodne są ciągłe, więc ograniczone.
Dla każdego wielomianu $P:\R^d\to\K$ funkcja
$f(x)=P(x)e^{-|x|^2}$ także należy do $\mathcal S(\R^d;\K)$.
Istotnie, $\partial_j(Pe^{-|x|^2})=
(\partial_jP-2x_jP)e^{-|x|^2}$, więc przez indukcję
każda pochodna ma postać wielomianu razy $e^{-|x|^2}$.
Po przemnożeniu przez dowolny jednomian otrzymujemy
funkcję ograniczoną, gdyż $(1+r)^k e^{-r^2}\to0$
dla każdego $k\in\N_0$ przy $r\to\infty$.
Natomiast $f(x)=(1+|x|^2)^{-1}$ nie należy do $\mathcal S$:
dla $t\ge1$ mamy $t^3|f(te_1)|=t^3/(1+t^2)\ge t/2$,
zatem $s_{3e_1,0}(f)=+\infty$.
\end{przyklad}

\begin{stwierdzenie}[Różniczkowanie, wielomiany i normy $L^p$]
\label{prop:schwartz-operators17}
Dla każdego $\alpha\in\N_0^d$ operatory
$f\mapsto\partial^\alpha f$ oraz $f\mapsto x^\alpha f$
są ciągłymi operatorami liniowymi na $\mathcal S(\R^d;\K)$.
Mnożenie $\mathcal S\times\mathcal S\to\mathcal S$ jest
łącznie ciągłe. Dla każdego $1\le p\le\infty$
kanoniczne włożenie $\mathcal S(\R^d;\K)\to L^p(\R^d;\K)$
jest ciągłe.
\end{stwierdzenie}
\begin{proof}
Różniczkowanie spełnia
$Q_{N,m}(\partial^\alpha f)\le Q_{N,m+|\alpha|}(f)$.
Dla mnożenia przez $x^\alpha$ wzór Leibniza daje
\[
\partial^\beta(x^\alpha f)=
\sum_{\substack{\gamma\le\beta\\\gamma\le\alpha}}
\binom{\beta}{\gamma}\frac{\alpha!}{(\alpha-\gamma)!}
x^{\alpha-\gamma}\partial^{\beta-\gamma}f.
\]
Ponieważ $|x^{\alpha-\gamma}|\le(1+|x|)^{|\alpha|}$,
po pomnożeniu przez wagę i wzięciu supremum otrzymujemy
$Q_{N,m}(x^\alpha f)\le C_{\alpha,m}Q_{N+|\alpha|,m}(f)$,
gdzie stała jest maksimum sum współczynników powyższego wzoru
dla $|\beta|\le m$. Analogiczny wzór daje
$Q_{N,m}(fg)\le C_{d,m}Q_{N,m}(f)Q_{0,m}(g)$,
a więc łączną ciągłość iloczynu.
Dla $p=\infty$ mamy $\|f\|_\infty\le Q_{0,0}(f)$.
Dla $1\le p<\infty$ wybieramy liczbę całkowitą $N>d/p$.
Wtedy
$\|f\|_{L^p}^p\le Q_{N,0}(f)^p
\int_{\R^d}(1+|x|)^{-Np}\,\mathrm d\lambda_d(x)$,
a całka jest skończona: na kuli jednostkowej całkujemy
funkcję ograniczoną, zaś w pierścieniu
$2^k\le|x|<2^{k+1}$, $k\ge0$, całka jest nie większa niż
$\lambda_d(B(0,1))2^d\,2^{k(d-Np)}$.
Szereg tych oszacowań jest zbieżny, bo $Np>d$.
Po wyciągnięciu pierwiastka stopnia $p$ dostajemy
żądane oszacowanie normy, a więc ciągłość włożenia.
\end{proof}

\begin{twierdzenie}[Gęstość funkcji o zwartym nośniku w $\mathcal S$]
\label{thm:density-schwartz17}
Przestrzeń $C_c^\infty(\R^d;\K)$ jest gęsta w
$\mathcal S(\R^d;\K)$ dla standardowej topologii Schwartza.
\end{twierdzenie}
\begin{proof}
Ustalmy $f\in\mathcal S(\R^d;\K)$.
Z Lematu~\ref{lem:cutoffs17} dla $K=\overline B(0,1)$
i $\Omega=B(0,2)$ wybieramy funkcję odcięcia, którą
przedłużamy zerem na $\R^d$ (jej nośnik jest zwarto zawarty
w $B(0,2)$, więc przedłużenie jest gładkie).
Otrzymujemy $\chi\in C_c^\infty(\R^d;\R)$ taką, że
$0\le\chi\le1$, $\chi=1$ na $\overline B(0,1)$
i $\operatorname{supp}\chi\subseteq B(0,2)$.
Niech $\chi_R(x)=\chi(x/R)$ dla $R\ge1$.
Funkcja $\chi_R f$ ma zwarty nośnik. Dla wielowskaźnika $\beta$
ze wzoru Leibniza
\[
\partial^\beta((1-\chi_R)f)
=(1-\chi_R)\partial^\beta f+
\sum_{0<\gamma\le\beta}\binom{\beta}{\gamma}
 (-1)R^{-|\gamma|}(\partial^\gamma\chi)(x/R)
 \partial^{\beta-\gamma}f.
\]
Nośnik pierwszego składnika jest zawarty w $\{|x|\ge R\}$,
a nośniki pozostałych w $\{R\le|x|\le2R\}$.
Dla $|\beta|\le m$ na nośniku każdego składnika mamy
$(1+|x|)^{-1}\le(1+R)^{-1}$. Ponadto
$|1-\chi_R|\le1$ i $R^{-|\gamma|}\le1$.
Po pomnożeniu przez $(1+|x|)^N$ pierwszy składnik jest
ograniczony przez $(1+R)^{-1}Q_{N+1,m}(f)$, a składnik
sumy o indeksie $\gamma$ przez tę samą wielkość
pomnożoną przez $\binom{\beta}{\gamma}\|\partial^\gamma\chi\|_\infty$.
Sumując, otrzymujemy
$Q_{N,m}((1-\chi_R)f)\le C_m(1+R)^{-1}Q_{N+1,m}(f)$,
gdzie $C_m$ jest maksimum odpowiednich skończonych sum
dla $|\beta|\le m$ i nie zależy od $R$ ani $f$.
Stąd $Q_{N,m}(f-\chi_Rf)\to0$ przy $R\to\infty$.
\end{proof}

\begin{stwierdzenie}[Brak normowalności]
\label{prop:non-normable17}
Jeśli $d\ge1$ i $\Omega\subseteq\R^d$ jest niepusty i otwarty,
to ani $C^\infty(\Omega;\K)$, ani $\mathcal S(\R^d;\K)$
ze swoimi standardowymi topologiami nie są normowalne.
\end{stwierdzenie}
\begin{proof}
Wystarczy pokazać, że żadne otoczenie zera nie jest ograniczone,
a następnie zastosować kryterium normowalności Kolmogorowa.
Najpierw rozważmy funkcje zespolone. Każde otoczenie zera
w $C^\infty(\Omega)$ zawiera zbiór $U=\{f:q_N(f)<\varepsilon\}$
dla pewnych $N\ge1$ i $\varepsilon>0$.
Wybierzmy kulę domkniętą o dodatnim promieniu zawartą w $\Omega$
i, z Lematu~\ref{lem:cutoffs17}, funkcję
$\psi\in C_c^\infty(\Omega;\R)$ równą $1$ w jej otoczeniu.
Jej nośnik zawiera się w pewnym $K_M$ z $M>N+1$.
Niech $f_j(x)=c j^{-N}\psi(x)e^{ijx_1}$.
Wzór Leibniza zapewnia $q_N(f_j)\le Cc$,
więc małym $c>0$ można zagwarantować $f_j\in U$.
Na wnętrzu wybranej kuli wszystkie pochodne $\psi$
dodatniego rzędu znikają, zatem
$\partial_1^{N+1}f_j(x)=c j i^{N+1}e^{ijx_1}$,
której moduł jest równy $cj$. Stąd
$q_M(f_j)\to\infty$, więc $U$ nie jest ograniczony.
Dla przestrzeni Schwartza wybieramy analogiczną funkcję
$\psi\in C_c^\infty(\R^d;\R)$, równą $1$ na otoczeniu kuli.
Na jej zwartym nośniku waga $(1+|x|)^N$ jest ograniczona,
więc z oszacowania Leibniza otrzymujemy $Q_N(f_j)\le C_Nc$.
Na wnętrzu kuli mamy natomiast
$Q_{N+1}(f_j)\ge cj$, ponieważ waga jest nie mniejsza od $1$.
Dla dowolnego otoczenia zera zawierającego
$\{f:Q_N(f)<\varepsilon\}$ wybieramy zatem
$0<c<\varepsilon/C_N$, uzyskując nieograniczony podzbiór
tego otoczenia.
Przypadek rzeczywisty otrzymujemy, stosując
$\cos(jx_1)$ zamiast $e^{ijx_1}$ i supremum na małej kuli,
gdzie $\psi$ jest stała i niezerowa; przy dużym $j$
cosinus lub sinus odpowiedniej pochodnej przyjmuje tam
wartość bezwzględną $1$.
\end{proof}

\subsection{Transformacja Fouriera w przestrzeni Schwartza}

W tej podsekcji pracujemy nad $\C$. Wzór z całką jest zwykłą
całką Lebesgue'a, bez potrzeby odwoływania się do dystrybucji.

\begin{definicja}[Transformacja Fouriera]
\label{def:fourier-schwartz17}
Dla $f\in\mathcal S(\R^d;\C)$ definiujemy
\[
\widehat f(\xi)=\mathcal Ff(\xi)=
(2\pi)^{-d/2}\int_{\R^d}e^{-ix\cdot\xi}f(x)\,
\mathrm d\lambda_d(x),\qquad\xi\in\R^d.
\]
Całka jest bezwzględnie zbieżna, gdyż $\mathcal S\subseteq L^1$.
\end{definicja}

\begin{twierdzenie}[Transformacja Fouriera zachowuje $\mathcal S$]
\label{thm:fourier-schwartz17}
Dla każdego $d\ge1$ transformacja Fouriera jest ciągłym
izomorfizmem liniowo-topologicznym
$\mathcal F:\mathcal S(\R^d;\C)\to\mathcal S(\R^d;\C)$.
Jej odwrotność wyraża wzór
\[
\mathcal F^{-1}g(x)=(2\pi)^{-d/2}
\int_{\R^d}e^{ix\cdot\xi}g(\xi)\,
\mathrm d\lambda_d(\xi).
\]
Dla każdego $f\in\mathcal S$ oraz $j=1,\dots,d$ zachodzą
$\mathcal F(\partial_jf)(\xi)=i\xi_j\widehat f(\xi)$
i $\mathcal F(x_jf)(\xi)=i\partial_{\xi_j}\widehat f(\xi)$.
\end{twierdzenie}
\begin{proof}
Przy różniczkowaniu względem $\xi_j$ iloraz różnicowy
czynnika wykładniczego jest ograniczony przez $|x_j|$.
Indukcyjnie majorantą przy kolejnych różniczkowaniach
jest $|x^{\beta+e_j}f(x)|$, całkowalna na mocy
Stwierdzenia~\ref{prop:schwartz-operators17}.
Twierdzenie o zbieżności zmajoryzowanej daje więc
$\partial_\xi^\beta\widehat f=(-i)^{|\beta|}
\mathcal F(x^\beta f)$, a także ciągłość tych pochodnych.
Dla pochodnej $\partial_j f$ stosujemy twierdzenie Fubiniego
i całkowanie przez części względem $x_j$ na $[-R,R]$.
Przy ustalonych pozostałych współrzędnych wyrazy brzegowe
dążą do zera, gdyż $f$ szybko maleje; całki jednowymiarowe
są bezwzględnie zbieżne. Następnie całkujemy po pozostałych
współrzędnych, korzystając z całkowalności $f$ i $\partial_j f$.
Otrzymujemy $\mathcal F(\partial_j f)=i\xi_j\widehat f$.
Iteracja daje $\mathcal F(\partial^\alpha f)=(i\xi)^\alpha\widehat f$.
Połączenie tych tożsamości prowadzi do
\[
\xi^\alpha\partial_\xi^\beta\widehat f(\xi)
=(-i)^{|\alpha|+|\beta|}
\mathcal F\bigl(\partial^\alpha(x^\beta f)\bigr)(\xi).
\]
Wartość bezwzględną prawej strony ograniczamy przez
$(2\pi)^{-d/2}\|\partial^\alpha(x^\beta f)\|_{L^1}$.
Z oszacowań w Stwierdzeniu~\ref{prop:schwartz-operators17}
otrzymujemy, dla pewnej stałej $C_{\alpha,\beta,d}$,
$s_{\alpha,\beta}(\widehat f)\le
C_{\alpha,\beta,d}Q_{d+1+|\beta|,|\alpha|}(f)$,
stosując kolejno oszacowanie normy $L^1$, różniczkowania
i mnożenia przez wielomian.
Wszystkie seminormy $s_{\alpha,\beta}(\widehat f)$ są więc
skończone i zależą w sposób ciągły od $f$.

Pozostaje udowodnić wzór odwrotny. Dla $\varepsilon>0$ definiujemy
\[
I_\varepsilon(x)=(2\pi)^{-d/2}\int_{\R^d}
 e^{-\varepsilon|\xi|^2/2}e^{ix\cdot\xi}
 \widehat f(\xi)\,\mathrm d\lambda_d(\xi).
\]
Ponieważ $f\in L^1$ i czynnik Gaussa należy do $L^1$,
możemy zastosować twierdzenie Fubiniego. Jednowymiarową całkę Gaussa obliczamy następująco:
$I(u)=\int_\R e^{-\varepsilon t^2/2}e^{iut}\,\mathrm dt$
można różniczkować pod całką z majorantą
$|t|e^{-\varepsilon t^2/2}$. Ponieważ
$(e^{-\varepsilon t^2/2})'=-\varepsilon t e^{-\varepsilon t^2/2}$,
całkowanie przez części, z zanikającymi wyrazami brzegowymi, daje
\[
I'(u)=-\frac{i}{\varepsilon}
\int_\R (e^{-\varepsilon t^2/2})'e^{iut}\,\mathrm dt
=-\frac u\varepsilon I(u).
\]
Ponadto $I(0)=\sqrt{2\pi/\varepsilon}$ z całki Gaussa
(po podstawieniu $s=\sqrt\varepsilon t$).
Zatem $I(u)=\sqrt{2\pi/\varepsilon}e^{-u^2/(2\varepsilon)}$.
Zastosowanie tego wzoru w każdej współrzędnej daje
\[
I_\varepsilon(x)=\int_{\R^d}f(y)
(2\pi\varepsilon)^{-d/2}e^{-|x-y|^2/(2\varepsilon)}
\,\mathrm d\lambda_d(y).
\]
Oznaczmy $G_\varepsilon(z)=(2\pi\varepsilon)^{-d/2}
e^{-|z|^2/(2\varepsilon)}$. Podstawienie $z=\sqrt\varepsilon w$
daje $\int G_\varepsilon\,\mathrm d\lambda_d=1$ oraz
$\int_{|z|\ge\delta}G_\varepsilon(z)\,\mathrm d\lambda_d(z)
=\int_{|w|\ge\delta/\sqrt\varepsilon}G_1(w)\,\mathrm d\lambda_d(w)
\to0$ dla każdego $\delta>0$, z całkowalności $G_1$.
Dla ustalonego $x$ i $\eta>0$ wybieramy $\delta>0$ tak,
aby $|f(x-z)-f(x)|<\eta$ dla $|z|<\delta$.
Wtedy
\[
|I_\varepsilon(x)-f(x)|\le
\eta+2\|f\|_\infty
\int_{|z|\ge\delta}G_\varepsilon(z)\,\mathrm d\lambda_d(z).
\]
Stąd granica górna błędu nie przekracza $\eta$, a ponieważ
$\eta$ jest dowolne, $I_\varepsilon(x)\to f(x)$.
Z drugiej strony $\widehat f\in L^1$, zatem
z twierdzenia o zbieżności zmajoryzowanej
$I_\varepsilon(x)$ dąży do całki ze wzoru odwrotnego.
Otrzymujemy $\mathcal F^{-1}\mathcal Ff=f$.
Ten sam argument, ze zmienionym znakiem wykładnika,
daje $\mathcal F\mathcal F^{-1}g=g$ dla $g\in\mathcal S$.
Ciągłość operatora odwrotnego wynika z identycznych
oszacowań seminorm po zmianie znaku wykładnika.
\end{proof}

\begin{przyklad}[Funkcja Gaussa i wizualizacja w Plotly]
\label{ex:fourier-gauss17}
Dla $a>0$ oraz $f_a(x)=e^{-ax^2/2}$, $x\in\R$,
stosując obliczoną w poprzednim dowodzie całkę $I(u)$
z $\varepsilon=a$ i $u=-\xi$, otrzymujemy
$\widehat f_a(\xi)=a^{-1/2}e^{-\xi^2/(2a)}$.
W szczególności zwiększenie skupienia funkcji $f_a$
odpowiada rozszerzeniu jej transformaty Fouriera.
Poniższy kod Python rysuje obie funkcje z użyciem biblioteki Plotly.
\begin{lstlisting}[language=Python,caption={Para funkcji Gaussa: funkcja i transformata Fouriera}]
import numpy as np
import plotly.graph_objects as go

a = 2.0
x = np.linspace(-5.0, 5.0, 1000)
f = np.exp(-a * x**2 / 2)
f_hat = a**(-0.5) * np.exp(-x**2 / (2 * a))

fig = go.Figure()
fig.add_scatter(x=x, y=f, mode="lines", name="f(x)")
fig.add_scatter(x=x, y=f_hat, mode="lines", name="Fourier(f)(x)")
fig.update_layout(title="Para Gaussa", xaxis_title="x / xi",
                  yaxis_title="wartosc")
fig.show()
\end{lstlisting}
\end{przyklad}

\begin{stwierdzenie}[Splot funkcji szybko malejących]
\label{prop:convolution-schwartz17}
Dla $f,g\in\mathcal S(\R^d;\C)$ definiujemy
$(f*g)(x)=\int_{\R^d}f(y)g(x-y)\,\mathrm d\lambda_d(y)$.
Wówczas $f*g\in\mathcal S$, odwzorowanie
$(f,g)\mapsto f*g$ jest dwuliniowe i łącznie ciągłe
na $\mathcal S\times\mathcal S$, a ponadto
$\mathcal F(f*g)=(2\pi)^{d/2}\widehat f\,\widehat g$.
\end{stwierdzenie}
\begin{proof}
Z $\mathcal S\subset L^1\cap L^\infty$ wynika
bezwzględna zbieżność splotu, gdyż całka modułu
nie przekracza $\|f\|_1\|g\|_\infty$.
Ilorazy różnicowe pochodnych $g$ względem $x_j$
są ograniczone przez $\|\partial^{\beta+e_j}g\|_\infty$,
więc po pomnożeniu przez $|f(y)|$ otrzymujemy majorantę
całkowalną względem $y$. Indukcja i twierdzenie o zbieżności
zmajoryzowanej dają $\partial^\beta(f*g)=f*(\partial^\beta g)$
i ciągłość wszystkich tych pochodnych.
Nierówność $1+|x|\le(1+|y|)(1+|x-y|)$ prowadzi do
\[
 Q_{N,m}(f*g)\le Q_{N,m}(g)
 \int_{\R^d}(1+|y|)^N|f(y)|\,\mathrm d\lambda_d(y)
 \le C_{d,N}Q_{N+d+1,0}(f)Q_{N,m}(g).
\]
To dowodzi jednocześnie szybkiego zanikania i ciągłości
splotu. Twierdzenie Tonellego i niezmienniczość miary Lebesgue'a
względem przesunięć dają
$\int_{\R^d}\int_{\R^d}|f(y)g(x-y)|\,\mathrm d\lambda_d(x)
\,\mathrm d\lambda_d(y)=\|f\|_1\|g\|_1<\infty$.
Możemy więc zastosować twierdzenie Fubiniego i podstawić $z=x-y$:
\[
\begin{aligned}
\mathcal F(f*g)(\xi)
&=(2\pi)^{-d/2}
\int_{\R^d} f(y)e^{-iy\cdot\xi}\,\mathrm d\lambda_d(y)
\int_{\R^d}g(z)e^{-iz\cdot\xi}\,\mathrm d\lambda_d(z)\\
&=(2\pi)^{d/2}\widehat f(\xi)\widehat g(\xi).
\end{aligned}
\]
\end{proof}

\begin{uwaga}[Przygotowanie teorii dystrybucji]
\label{rem:distribution-next17}
Ciągłe funkcjonały na $C^\infty(\Omega)$ mają oszacowania
na pochodne na jednym zwartym zbiorze (Twierdzenie~\ref{thm:dual-cinf17}).
Dla liniowego funkcjonału $T:\mathcal S(\R^d;\K)\to\K$
analogiczny warunek ciągłości ma postać
$|T(f)|\le C Q_{N,m}(f)$ dla pewnych $C>0$, $N,m\in\N_0$
i wszystkich $f\in\mathcal S(\R^d;\K)$.
Istotnie, przy ciągłości pewne otoczenie
$\{f:Q_n(f)<\varepsilon\}$ jest odwzorowywane w kulę jednostkową
ciała skalarów. Skalowanie, jak w dowodzie
Twierdzenia~\ref{thm:dual-cinf17}, daje
$|T(f)|\le\varepsilon^{-1}Q_{n,n}(f)$.
Odwrotnie, takie oszacowanie zapewnia ciągłość w zerze,
a z liniowości również w każdym punkcie.
Przestrzenie ciągłych funkcjonałów liniowych na tych
przestrzeniach funkcji nie są tożsame.
W następnym rozdziale skonstruujemy $\mathcal D(\Omega)$
z topologią granicy induktywnej przestrzeni funkcji o nośniku
zawartym w ustalonym zbiorze zwartym. Dopiero wtedy zdefiniujemy
ogólną przestrzeń dystrybucji $\mathcal D'(\Omega)$.
\end{uwaga}
